% Part: first-order-logic
% Chapter: sequent-calculus

\documentclass[../../../include/open-logic-chapter]{subfiles}

\begin{document}

\iftag{FOL}
      {\olchapter{fol}{seq}{સિક્વન્ટ કલન}}
      {\olchapter{pl}{seq}{સિક્વન્ટ કલન}}

\begin{editorial}
  આ પ્રકરણ શાસ્ત્રીય પ્રથમ-ક્રમના તર્કશાસ્ત્ર માટે ગેન્ટ્ઝનનું પ્રમાણભૂત
  સિક્વન્ટ કલન LK રજૂ કરે છે. તેમાં હજી વધુ દાખલાઓ અને
  અભ્યાસપ્રશ્નો ઉમેરવા યોગ્ય છે. સાબિતી-તંત્ર તરીકે સિક્વન્ટ કલનને લાગતી
  સામગ્રીનો સમાવેશ કરવા કે તેને બાકાત રાખવા માટે ``prfSC'' ટૅગ વાપરો.
\sourcecorrection{OLREA-235}{મૂળની \texttt{sequent-calculus.tex}, પંક્તિ 16માં સંપાદકીય સૂચના \texttt{prfLK} ટૅગ આપે છે. સ્થિર મૂળની \texttt{open-logic-config.sty}, પંક્તિઓ 1629--1637 અને બંને મુખ્ય તર્કશાસ્ત્ર-ભાગોના આયાત-નિયમોમાં સિક્વન્ટ કલનનો નિયામક ટૅગ \texttt{prfSC} છે. અહીં વપરાતા ટૅગનું સાચું નામ આપ્યું છે.}
આ આવૃત્તિમાં આ પ્રકરણનું $!A\liff!B$ એ
$(!A\lif!B)\land(!B\lif!A)$નું સંક્ષિપ્ત સ્વરૂપ છે; નિયમો લાગુ
કરતાં પહેલાં તેને વિસ્તારી લઈએ. $\ltrue$ અને $\lfalse$ માટે
આગળની પ્રારંભિક સિક્વન્ટોની વ્યાખ્યામાં જ નિયમો આપેલા છે.
\end{editorial}

\olimport{rules-and-proofs}

\olimport{propositional-rules}

\iftag{FOL}{%
  \olimport{quantifier-rules}
}{}

\olimport{structural-rules}

\olimport{derivations}

\olimport{proving-things}

\iftag{FOL}{%
  \olimport{proving-things-quant}
}{}

\olimport{proof-theoretic-notions}

\olimport{provability-consistency}

\olimport{provability-propositional}

\iftag{FOL}{%
  \olimport{provability-quantifiers}
}{}

\olimport{soundness}

\iftag{FOL}{%
  \olimport{identity}
  \olimport{soundness-identity}
}{}

\OLEndChapterHook

\end{document}
